\section{Conclusions}
\label{sec:summary}

Starting with a Euclidean operator product expansion for products of gauge invariant operators in QCD, we have derived the factorization formulas for the quasi-PDF, \spcorr and pseudo-PDF. The three Euclidean distribution functions are related observables, and all follow from the same fundamental factorization.  For the \spcorr this derivation implies that the ratio in Eq.~(\ref{eq:ratio}) does scale in $z^2$, but needs the small $z^2$ factorization formula in \eq{io-Q-fact} to extract the PDF.
%
Our derivation for the factorization formula applies when the renormalized \spcorr is defined in any scheme.
The OPE used here could also be used to systematically derive factorization formulas for power corrections to \eq{momfact}, which will involve matching onto higher-twist parton distributions. (The numerical relevance of these corrections is considered in Ref.~\cite{Chen:2016utp}.)
%
Note that LaMET is not equivalent to the expansion from the OPE, as the former is more general and can be applied to the lattice calculations of other quantities, for example the TMD-PDF where a simple OPE does not exist.

Our derivation of the factorization formula for the quasi-PDF also verifies that the parton momentum $p^z$ in the matching coefficient in Eq.~(\ref{eq:momfact}) has to be $p^z=yP^z$, which makes a considerable difference for the $\overline{\rm MS}$  matching result when compared with $p^z=P^z$ (see Fig.~\ref{fig:quasi}). The proper $p^z$ should therefore be used in lattice calculations of the PDF in the LaMET approach.

As a non-trivial test of relations between the various distributions and factorization formulas we have considered results at one-loop in the $\overline{\rm MS}$ scheme. We have derived a corrected results for the coefficient $C$ for the $\overline{\rm MS}$ scheme, given in \eq{quasi-matching}, which leads to convergent results in the convolution integral.   We have also computed the one-loop $\overline{\rm MS}$ result for the Wilson coefficient ${\cal C}$ appearing in the \spcorr and pseudo-PDF factorizations. A numerical analysis of these one-loop corrections in $\overline{\rm MS}$ has also been provided. The one-loop matching coefficient ${\cal C}$ has a smaller effect for the pseudo-PDF than $C$ does for quasi-PDF, as can seen by comparing \figs{quasi}{pseudo}. Given systematic uncertainties in manipulating the lattice data, it is potentially interesting to consider using the same lattice data on the \spcorr to extract the parton distribution function using both the quasi-PDF and pseudo-PDF approaches.

There are several different ways of implementing the factorization formula to calculate the PDF from lattice data for the \spcorr $\tilde Q$, which we have discussed in \sec{lattice}. One always has to work with short distance correlation and large nucleon momentum to reduce higher-twist corrections. This limits the number of useful data points from lattice calculations as described in \sec{lattice}.  To achieve precision calculations without making model assumptions it will be highly desirable to move towards finer lattice spacing to increase the number of effective data points.
